% !TeX root = ../papers/nc-p3.tex
\section{有理正規三次の割線と四直線対称行列網}
\label{sec:rank-one-twisted-cubic-exclusion}
\begin{theorem}\label{thm:rank-one-twisted-cubic-exclusion}
被約な四辺有限の標準有限Frobenius窓の、一般可逆で全二次式を成分に持つ反対根は、
階数一軌跡に既約な三次曲線を持たない。
\end{theorem}
\begin{proof}[交代六座標が強制する行列網]
三次成分は平面的ではあり得ず、有理正規三次 $C$ である。反対側も平面的ではない。
実際完全な二次写像の包の次元は $C$ 上7で、平面では6以下である。
両側とも完全な三次系なので座標を合わせると $N|_C=xx^t$ となる。三つの標準二次式の列 $q(x)$ と
対称一次行列の網 $A$ を用いて $N(x)=xx^t+A(q(x))$ と書く。
パラメータ $t=[q(x)]$ は $C$ の長さ二因子の平面で、$C$ の外のその繊維は対応する割線 $K_t$ である。
実際三つの $q$ は二行三列Hankel行列の符号付き二行小行列式で、その核 $t$ を固定した二式は
長さ二因子の包の直線を切る。接線では非被約な二因子を用いる。
$K_t$ 上の $q(x)$ はその因子で零となる二次式と定数列 $t$ の積である。

印で $A(q)=yy^t-xx^t$ であるから、六つの交代評価 $x\wedge y$ が非零のところは階数2で、
そのベクトルは像二平面のPlücker座標に比例する。十六の混合積が独立なので、このようなPlücker列は
六次元を張らなければならない。$A$ の一般階数2以下は $N$ の一般可逆性に反する。
$A$ の係数の定数核もない。共通零化列 $v$ があれば $Nv=(x^tv)x$ が平面因子を持つからである。
一般階数3なら三次余因子行列を $hgg^t$ と原始に書く。$g$ は非定数だから一次、$h$ は一次で、
$g$ の基点は高々一点となる。階数二以下のパラメータはこの点と直線 $h=0$ に入り、
全印は二次曲面 $h(q)=0$ と高々一本の $K_t$ に含まれる。$14+3<20$ で不可能である。
従って $\det A$ は非零四次式である。係数写像の像が二次元以下なら一般可逆な行列鉛筆の階数2点は
四次行列式の二重根なので高々二つ、零行列は交代評価に寄与せず、六つのPlücker方向を張れない。
\end{proof}
\begin{proof}[四次判別式の全分割]
既約四次判別曲線の特異点は高々三つで、被約な分割 $(3,1),(2,2),(2,1,1)$ の特異点もそれぞれ
高々 $1+3,4,1+2+2$ である。階数2点は特異なので、被約な場合に六方向を持つには四直線があり、
三直線が共点にならず、全六頂点で階数2となることが必要である。

非被約な場合を確かめる。重複成分が滑らかな二次曲線なら他の成分はない。
その上で一般階数2の場合、第二外冪の二項四次式を $hgg^t$ と書くとPlücker列 $g$ の次数は高々2で、
得られる方向の包は三次元以下。一般階数3なら三次余因子の二項六次式での同じ原始分解を用いる。
核一定なら二次曲線はパラメータ平面を張るので網全体の固定核となり不可能である。
核が動けばスカラー因子の次数は4以下で階数二点は四つ以下である。一般階数1なら寄与はない。

重複成分が直線なら、その上の一般階数2のときPlücker列の次数は高々1で寄与次元は2以下、
一般階数3なら余因子の次数から階数二点は三つ以下で、三つあれば核は直線上で一定である。
その直線の外の階数二点は残余二次の特異点なので、別の重複直線がない限り高々一つである。
二直線がともに二重なら六方向には各々で三点が必要で、各直線の核は一定でなければならない。
核が同じなら網全体の核となり、異なれば交点が既に階数2以下なので二組に重複が生じる。
これで非被約な判別式を全て除いた。

残った四直線 $h_i(t)=0$ の各々で余因子が三頂点を通るから、やはり核は一定である。
核列 $v_i$ には $A(t)v_i=h_i(t)u_i$ が成り立つ。
一般の二点 $t_0,t_1$ を選べば $A(t_0)^{-1}A(t_1)$ の四固有値
$h_i(t_1)/h_i(t_0)$ は相異なり、四核列は独立となる。対称性でこの基底は $A(t_0)$ に直交し、
$A$ は定数合同で対角となる。$x$ の座標も合わせ、以下対角成分を定数倍込みで $h_i$ とする。
\end{proof}
\begin{proof}[四面体の辺と印の上界]
階数二パラメータ $t_{ij}$ を、$i,j$ の二対角成分だけが非零となる頂点とする。
その像は座標辺 $E_{ij}=\langle e_i,e_j\rangle$ である。曲線外の印は
$K_{t_{ij}}\cap E_{ij}$ に入り、全六辺で少なくとも一つの曲線外印が要る。実際交代評価の各座標はその辺だけが与える。
二直線が異なれば一点以下である。一致する辺を選んで四頂点グラフを作る。一頂点に三つの選ばれた辺は集まらない。
異なる割線が交わるとその点は $C$ 上であり、対応する二因子がその点を共有する。
これは有理正規三次で長さ4以下の因子が独立な線型条件を持つためである。
一頂点での三因子は従って因子平面内の一直線に並ぶはずだが、
$t_{ij}$ を $i$ 固定、$j$ 可変で取ると他の三直線の交点からなる三角形になる。
ゆえに頂点次数は2以下。もし四辺以上ならちょうど四角形で各頂点は $C$ 上にある。
選ばれなかった対角辺も二頂点の割線なので、異なる $K_{t_{ij}}$ とは $C$ の外で交われない。
必要な交代評価に反し、選ばれた辺は高々三本である。

一致辺上で $x=(a,b)$、対応する $A(q)$ の非零部分を
$\lambda(a,b)\operatorname{diag}(c_i,c_j)$ と書く。$\lambda$ の零因子は $C$ との長さ二交点である。
二行行列式は
\[
       \lambda\big(c_j a^2+c_i b^2+c_ic_j\lambda\big).
\]
よって辺全体が階数一でないとき、$C$ 外の印は二つ以下である。辺全体が階数一なら $\lambda$ は
両座標頂点で零でない。従ってその辺と頂点を共有する別の一致辺はない。
全体が階数一の辺がなければ曲線外は $6+3=9$ 印以下だが、$C$ の上は三次評価から十印以下である。
一本全体が階数一なら、別に一致できるのは反対の辺だけである。その反対辺の全体が階数一でなければ
曲線外は $3+2+4=9$ 印以下で再び不可能となる。
最後に二本の反対辺がともに階数一なら、$C$ とこの二本は連結な総次数5の曲線で、残る四辺上にはその外の印が高々四つである。
この曲線上で混合束 $LL'$ の切断は $2\cdot5+1=11$ 次元以下
（定理\ref{thm:rank-one-curve-degree-bound}の連結曲線の上界）。
従って全印でも混合積は $11+4<16$ 次元となる。これで最後の場合も除かれた。
\end{proof}
